\chapter{Complementos da fundamentação teórica}
\label{ap:fundamentacao}

Este apêndice reúne o detalhamento dos modelos de transformação digital, das opções de arquitetura e dos construtos de aceitação usados na fundamentação.

\section{Transformação digital e pequenos negócios}

Para avaliar em que estágio desse processo um estabelecimento está, \citeonline{westerman2014leading} propõem um modelo de maturidade digital, baseado em um estudo do MIT com a Capgemini. O modelo tem duas dimensões independentes. A primeira é a capacidade digital, ou seja, quanto a empresa investe em tecnologias que mudam a forma de operar. A segunda é a capacidade de liderança, que envolve o apoio da direção, a governança e o engajamento das pessoas na transformação. O cruzamento das duas dimensões gera quatro quadrantes, \emph{Beginners} (baixa capacidade nas duas), \emph{Fashionistas} (alta capacidade digital, mas baixa capacidade de liderança), \emph{Conservatives} (situação inversa) e \emph{Digital Masters} (alta capacidade nas duas), com diferenças consistentes de desempenho financeiro entre eles. Esse modelo ajuda a interpretar as diferenças entre os perfis estudados. Um restaurante que adota cardápio digital sem rever processos nem treinar a equipe (\emph{Fashionista}) tende a ter resultados diferentes de outro que liga a tecnologia a uma visão clara de atendimento (\emph{Digital Master}). Isso pode acontecer mesmo que os dois usem a mesma plataforma.

Em estudo com sete pequenos negócios alimentícios no estado de Sergipe, \citeonline{costa2021implicacoes} demonstram, por meio de entrevistas semiestruturadas analisadas tematicamente, que a pandemia atuou como catalisador externo da inserção de tecnologia, frequentemente sem que houvesse aceitação prévia dos proprietários. \citeonline{harms2021effectuation}, em pesquisa com 143 empreendedores gastronômicos em Münster, Alemanha, identificam dois caminhos de inovação de modelo de negócio sob alta incerteza: o do \emph{planning soloist}, baseado em causation (planejamento sequencial), e o do \emph{hedging networker}, baseado em effectuation (uso de recursos disponíveis e parcerias). Ambos os caminhos culminam, com frequência, em ofertas digitais coordenadas, como o ``\emph{digital beer garden}'' cooperativo descrito pelos autores.

Estudos mais recentes indicam que a digitalização desencadeada pela pandemia tende a ser estrutural, e não conjuntural. \citeonline{giovanini2024pandemia}, por meio de simulação computacional com modelos baseados em agentes, mostram que a difusão de plataformas digitais de entrega sob demanda persiste mesmo após o fim do choque pandêmico, em razão de mecanismos de feedback positivo entre consumidores e restaurantes. No campo das culturas alimentares mais amplas, \citeonline{lupton2020digitalfood} sintetizam, em coletânea publicada pela Routledge, como mídias sociais, dispositivos de self-tracking e interfaces digitais estão remodelando práticas, identidades e poder simbólico em torno do alimento. Para este trabalho, isso tem três consequências. A digitalização do food service brasileiro tem um marco claro antes e depois de 2020. Em pequenas e médias empresas, as decisões sobre tecnologia são tomadas com incerteza e pouco capital. E o cardápio digital, por estar visível no momento do consumo, é um bom ponto para observar essa transformação na relação com o cliente.

\section{Plataformas e opções de arquitetura}

O modelo SaaS caracteriza-se pela entrega de software como serviço acessado via navegador, com atualizações contínuas, dispensando aquisição de licenças e gestão de infraestrutura local pelo cliente \cite{correia2024crescimentosaas}. Em uma revisão sobre o crescimento do SaaS, \citeonline{correia2024crescimentosaas} apontam como principais fatores a flexibilidade da assinatura, o acesso mais fácil a tecnologias atuais e a adaptação a diferentes áreas. Por outro lado, os autores destacam que segurança da informação e conformidade legal ainda são obstáculos, principalmente quando há dados de clientes envolvidos, ponto retomado na Seção~\ref{sec:lgpd}.

Do ponto de vista arquitetural, plataformas SaaS de cardápio digital podem operar em regime multi-tenant, em que uma única instância da aplicação serve simultaneamente múltiplos clientes (tenants), com isolamento lógico de dados e configurações. \citeonline{rodriguesneto2012saas} sistematizam os principais padrões para implementação de aplicações multi-tenant, discutindo as três abordagens clássicas de isolamento de dados (banco de dados separado por tenant, esquema separado em banco compartilhado, ou tabelas compartilhadas com discriminador de tenant) e seus trade-offs entre custo de infraestrutura, complexidade operacional e grau de personalização possível. Essa escolha afeta diretamente o fornecedor de cardápio digital. O modelo de tabelas compartilhadas aproveita melhor a infraestrutura e reduz o custo de cada novo estabelecimento, o que torna viável atender muitos pequenos negócios. Já o modelo de banco separado pode ser adotado quando houver exigências maiores de isolamento dos dados.

Para o gestor do estabelecimento, a contratação de um SaaS reduz a necessidade de manter equipe técnica própria, mas introduz dependência do fornecedor, cuja estabilidade comercial e técnica passa a integrar o risco operacional do negócio. A entrevista com o fornecedor (Apêndice~\ref{ap:roteiro-fornecedor}) busca verificar, sem presumir a arquitetura das plataformas observadas, como se dão hospedagem, isolamento de dados e personalização.

\textbf{Camadas e pontos de integração.} Além das três camadas apresentadas na Seção~\ref{sec:cardapios-digitais}, um self-ordering completo tem dois pontos críticos de integração. O primeiro é com o PDV/ERP do estabelecimento, para registrar a comanda, atualizar o estoque e emitir o documento fiscal. O segundo é com a cozinha, por impressora térmica, tablet de produção ou Kitchen Display System (KDS). A qualidade dessas integrações define, em larga medida, o ganho operacional percebido pelo gestor: um cardápio digital sem integração com PDV gera retrabalho de digitação manual; sem integração com a cozinha, perde o efeito de redução de tempo de atendimento que justifica a adoção \cite{fioreti2026aplicativos,oliveira2017autoatendimento}.

\textbf{Conectividade e estratégias de tolerância a falhas.} Como destacado em \citeonline{bass2021software}, disponibilidade não é uma propriedade gratuita: depende de redundância, monitoramento e estratégias explícitas de degradação graciosa. No food service, há duas frentes de fragilidade. A primeira é a conectividade do estabelecimento: se a internet do salão cai, o back-end torna-se inacessível e o fluxo de pedido se interrompe. Estratégias como cache local no tablet de cozinha e modos offline-first para o front-end mitigam essa fragilidade. A segunda é a conectividade do dispositivo do cliente: como o aparelho usado é do próprio consumidor, o estabelecimento não tem controle sobre cobertura de rede móvel ou consumo de dados; a oferta de Wi-Fi gratuito tornou-se, por isso, requisito de fato para a viabilidade do autoatendimento, e, em Porto Alegre, requisito normativo, conforme a lei municipal mencionada na Seção~\ref{sec:lgpd}.

\textbf{Disponibilidade como requisito não funcional crítico.} A indisponibilidade do cardápio digital durante o pico de movimento é, na prática, uma falha grave: o cliente fica sem acesso ao menu e o garçom precisa improvisar atendimento manual. Em termos da norma \cite{iso25010}, isso corresponde às subcaracterísticas de \emph{maturity}, \emph{availability} e \emph{recoverability} dentro da característica de confiabilidade. A literatura empírica sobre incidentes e SLAs efetivos dessas plataformas é ainda escassa, configurando uma lacuna que estudos empíricos com gestores podem ajudar a iluminar.

\section{Usabilidade e inclusão no autoatendimento}

\citeonline{oliveira2017autoatendimento}, em pesquisa com clientes de uma rede de \emph{fast food} em São Paulo, identificaram percepção positiva quanto a conforto e rapidez associados ao autoatendimento, mas observaram que esse fator não foi decisivo na escolha do estabelecimento. \citeonline{fioreti2026aplicativos} estudaram 100 clientes e alguns estabelecimentos do setor, com métodos qualitativos e quantitativos. Entre os clientes, 80\% mostraram alta intenção de usar aplicativos de autoatendimento, citando conveniência (85\%), rapidez (78\%) e mais controle sobre o pedido (72\%). Do lado dos estabelecimentos, os autores relatam redução de até 30\% no tempo médio de atendimento e aumento de até 20\% no tíquete médio. Em paralelo, os autores destacam desafios persistentes relativos à usabilidade das interfaces, à familiaridade tecnológica de parte dos usuários e à necessidade de treinamento das equipes.

No domínio específico de pedidos por aplicativos, \citeonline{martiniano2022delivery}, ao realizarem 12 entrevistas em profundidade com usuários brasileiros de \emph{Mobile Food Ordering Apps} (MFOAs), identificaram cinco famílias de determinantes de uso: influenciadores, valor percebido, disponibilidade do aplicativo, predileção por um app específico e dificuldade de uso. Observaram também que as avaliações dentro do app (sobre o restaurante) são mais relevantes para a decisão do que as avaliações globais do app na loja, e que a motivação hedônica e o estado emocional do consumidor influenciam diretamente o momento da compra. Esses achados, embora coletados em contexto de delivery, oferecem hipóteses transferíveis para o cenário presencial estudado neste trabalho.

\citeonline{mclaughlin2020designing} reúnem estudos de fatores humanos e de psicologia do envelhecimento aplicados ao projeto de telas. Os autores descrevem as limitações que costumam aparecer com a idade: visão menos nítida e menor sensibilidade ao contraste, menos precisão em movimentos finos, processamento mais lento e maior dependência de pistas visuais claras, em vez da memória de curto prazo. \citeonline{lara2010acessibilidade} estudaram adultos de meia-idade (40 a 59 anos) e idosos (60 anos ou mais) usando uma loja virtual. Recursos como aumentar a fonte, títulos de página claros, mensagens de erro ao lado do campo com problema e divisão da tarefa em passos menores reduziram a diferença de desempenho entre as faixas etárias. Esses recursos correspondem a critérios da WCAG e podem ser aplicados ao pedido por QR Code.

No plano do projeto de interfaces inclusivas, \citeonline{modolo2024designinclusivo} aplicam princípios de Design Universal \cite{mace1997principles} a aplicativos móveis voltados a usuários idosos, demonstrando que escolhas de colorimetria (paleta de cores tríade considerando condições de iluminação variáveis), tipografia (corpos amplos e fontes de alto contraste) e iconografia (ícones reconhecíveis com rótulos textuais) podem assegurar acessibilidade sem comprometer funcionalidade ou estética. Cardápios digitais com forte apelo visual (fotos de pratos, selos de promoção) têm justamente nessas dimensões seus pontos mais frágeis quando avaliados com base nesses princípios. Modalidades alternativas de interação, como interfaces conversacionais por chatbot, foram exploradas por \citeonline{oliveira2023interfaces} em estudo com usuários idosos: os achados sugerem aceitação positiva quando o design considera presença social e linguagem coloquial, mas a modalidade conversacional permanece tangencial ao recorte deste TCC, em que o QR Code remete a uma interface visual típica.

As dez heurísticas de \citeonline{nielsen1994heuristics}, complementadas por \citeonline{nielsen1993usability}, incluem visibilidade do status do sistema, correspondência com o mundo real, controle e liberdade do usuário, consistência, prevenção de erros, reconhecimento em vez de recordação, flexibilidade e eficiência, design minimalista, ajuda na recuperação de erros, e ajuda e documentação. No contexto de cardápios por QR Code, problemas mapeáveis a essas heurísticas incluem a ausência de feedback sobre envio do pedido, a falta de botão para voltar e mensagens de erro genéricas em caso de falha de rede.

\section{Modelos de aceitação e operacionalização}

\citeonline{davis1989useracceptance} compararam empiricamente o TAM com a TRA e mostraram que, embora o TAM seja mais específico ao domínio de tecnologia, sua capacidade preditiva é equivalente à da TRA com instrumentos mais enxutos. \citeonline{davis1993user} consolidou o modelo causal completo, conectando características do sistema (variáveis externas) às percepções de utilidade e facilidade, à atitude, à intenção e ao uso efetivo. \citeonline{king2006metaanalysis}, em meta-análise de 88 estudos empíricos, confirmaram a robustez e a validade do TAM, ressaltando, contudo, variação significativa dos efeitos em função do tipo de usuário e de sistema. \citeonline{legris2003why}, em revisão crítica anterior, já apontavam que TAM e sua extensão TAM2 explicavam aproximadamente 40\% da variância no uso efetivo, deixando uma fração substancial sem explicação.

Como resposta a essa fragmentação, \citeonline{venkatesh2003useracceptance} propuseram a Unified Theory of Acceptance and Use of Technology (UTAUT), unificando oito modelos prévios (TRA, TAM, Motivational Model, Theory of Planned Behavior, modelo combinado TAM-TPB, Model of PC Utilization, Innovation Diffusion Theory e Social Cognitive Theory) em um único modelo. A UTAUT identifica quatro construções nucleares, expectativa de performance, expectativa de esforço, influência social e condições facilitadoras, e quatro moderadores (idade, gênero, experiência e voluntariedade de uso). O modelo alcança R² ajustado de cerca de 69\% na previsão da intenção de uso, desempenho consideravelmente superior ao dos modelos anteriores tomados isoladamente.

A literatura brasileira tem incorporado a UTAUT2 em diferentes recortes. \citeonline{martiniano2022delivery} usam a UTAUT2 para discutir o comportamento de usuários de aplicativos de pedido de comida. \citeonline{matte2021evolucao}, em uma revisão de 145 artigos, mostram que a UTAUT2 está entre as teorias de adoção mais recentes e ainda pouco usadas no Brasil. Já \citeonline{gonzales2017teoria} mapearam os trabalhos brasileiros sobre a UTAUT em três eventos científicos (CONTECSI, ENANPAD e ENADI) entre 2011 e 2015 e observaram que eles se concentram em educação e comércio.

Para apoio didático, \citeonline{marikyan2024tam} e \citeonline{marikyan2023utaut}, no contexto do \emph{TheoryHub Book} da Newcastle University, fornecem verbetes de referência sobre o TAM e a UTAUT, respectivamente, com síntese histórica, aplicações e críticas contemporâneas.

\begin{table}[!ht]
    \caption{Construtos da UTAUT2 adotados na perspectiva do consumidor}
    \label{tab:utaut2-construtos}
    \centering%
    \begin{minipage}{.98\textwidth}
        \small
        \renewcommand{\arraystretch}{1.3}
        \begin{tabular*}{\linewidth}{@{}p{3.1cm}p{5.2cm}p{5.2cm}@{}}
            \toprule
            \textbf{Construto} & \textbf{Definição adotada} & \textbf{Exemplo no contexto do estudo} \\
            \midrule
            Expectativa de desempenho (PE) & Grau em que o cliente acredita que o cardápio digital melhora seu desempenho na tarefa de pedir & ``é mais prático que o cardápio físico''; ``ajuda a escolher mais rapidamente'' \\
            Expectativa de esforço (EE) & Grau de facilidade percebida no uso do cardápio digital & ``é fácil aprender a usar''; ``tenho dificuldade de navegar'' (invertido) \\
            Influência social (SI) & Grau em que o cliente percebe que outros importantes para ele acham que deveria usar o cardápio digital & ``pessoas próximas preferem estabelecimentos com cardápio digital''; ``avaliações em redes sociais influenciam'' \\
            Condições facilitadoras (FC) & Grau em que recursos e apoio necessários ao uso estão disponíveis para o cliente & ``tenho conexão de internet adequada''; ``há alguém para auxiliar quando preciso'' \\
            Motivação hedônica (HM) & Grau em que o uso do cardápio digital é percebido como prazeroso & ``acho prazeroso explorar com fotos e descrições''; ``atrapalha a experiência social'' (invertido) \\
            Valor (PV) & Avaliação do custo-benefício (tempo, dados móveis, indução de gasto) versus benefícios percebidos & ``vale a pena usar meus dados móveis''; ``me leva a gastar mais'' (invertido) \\
            Hábito (HB) & Grau em que o uso do cardápio digital se tornou automático & ``virou um hábito automático''; ``uso sem precisar pensar muito'' \\
            Intenção comportamental (BI) & Probabilidade declarada de continuar usando o cardápio digital & ``pretendo continuar usando nos próximos meses''; ``recomendaria a um amigo'' \\
            \bottomrule
        \end{tabular*}
        \fonte{Elaborado pelo autor a partir de \citeonline{venkatesh2012consumer}.}
    \end{minipage}
\end{table}
